\subsection{Determinants}

% TODO - Catch up on previous lecture cofactor reduction

\paragraph*{Theorem}
Let $A$ be an $n \times n$ matrix
\begin{enumerate}
    \item If A has a row or column of zeros, then $det(A) = 0$
    \item If two distinct rows or columns are interchanged, the determinant changes sign
    \item If a row is multiplied by a scalar $k$, then the determinant is multiplied by $k$
    \item $det(kA) = k^n det(A)$ and $det(-A) = -det(A)$
    \item If two rows or columns are identical or scalar multiples of each other, then $det(A) = 0$
    \item If a multiple of one row or column is added to another row or column, the determinant doesn't change
\end{enumerate}

\paragraph*{Example}
Find the determinant of \begin{vmatrix}
    1 & 1 & 1 & 2 \\
    2 & 2 & 0 & 1 \\
    3 & 3 & -1 & -3 \\
    4 & 2 & -5 & 7
\end{vmatrix}

\begin{align*}
    \text{det} &= \begin{vmatrix}
        1 & 1 & 1 & 2 \\
        2 & 2 & 0 & 1 \\
        3 & 3 & -1 & -3 \\
        4 & 2 & -5 & 7
    \end{vmatrix} \\
       &= \begin{vmatrix}
        1 & 1 & 1 & 2 \\
        0 & 2 & 0 & 1 \\
        0 & 3 & -1 & -3 \\
        -2 & 2 & -5 & 7
    \end{vmatrix}
    &= -2 \begin{vmatrix}
        1 & 1 & 2 \\
        2 & 0 & 1 \\
        3 & -1 & -3
    \end{vmatrix}
    &= -2 (-1) \begin{vmatrix}
        2 & 1 \\
        4 & -1
    \end{vmatrix}
    &= 2(-2-4) = -12
\end{align*}

\paragraph*{Example}
\det\begin{bmatrix}
    a & b & c \\
    p & q & r \\
    x & y & z
\end{bmatrix} = 6. Evauluate $\det\begin{bmatrix}
    a + x & b + y & c + z \\
    3x & 3y & 3z \\
    -p & -q & -r
\end{bmatrix}$
\begin{align*}
    = -3 (-1) \begin{vmatrix}
        a + x & b + y & c + z \\
        x & y & z\\
        p & q & r
    \end{vmatrix}
    = (-1)(-3) \begin{vmatrix}
        a & b & c \\
        p & q & r \\
        x & y & z
    \end{vmatrix} = 3 \cdot 6 = 18
\end{align*}

\subsubsection*{Determinants of Elementary Matricies}
The type 1 elementary matrix $E_1$ is obtained by interchanging two rows.
The type 2 elementary matrix $E_2$ is obtained by multiplying a row by a nonzero scalar $k$.
The type 3 elementary matrix $E_3$ is obtained by adding a multiple of one row to another row.
\begin{itemize}
    \item $\det(I) = 1$
    \item $\det(E_1) = -1$
    \item $\det(E_2) = k$
    \item $\det(E_3) = 1$
\end{itemize}

\subsubsection*{Determinants of Triangular Matrices}
\paragraph*{Definition}
\begin{itemize}
    \item $A_{n \times n}$ is called a lower triangular matrix if all entries above the main diagonal are zeros. 
    \item $A$ is called an upper triangular matrix if all entries below the main diagonal are zeros.
    \item $A$ is a triangular matrix if it is either an upper or lower triangular matrix.
    \item $A$ is a diagonal matrix if all entries above and below the main diagonal are zeros.
    \item If $A$ is diagnonal, then it is also triangular.
\end{itemize}

\paragraph*{Theorem}
If $A_{n \times n}$ is triangular, then $\det(A) = a_{11}a_{22} \cdots a_{nn}$.

\subsubsection*{The Geometric Meaning of the Determinant}
\paragraph*{1. $\mathbb{R}^2$}
Let $A_{n \times n}$ and $T_A(\vec{x}) = A\vec{x}$. In $\mathbb{R}^2$, $\det(A)$ gives the area-scaling factor of $T_A$
This means if $S$ is a region in $\mathbb{R}$, then $Area(T_A(S)) = |\det(A)| \cdot Area(S)$. The area of the parallelogram spanned by 
vectors $\vec{u}$ and $\vec{v}$ is given by $A = |\det([\vec{u}, \vec{v}])|$

\paragraph*{2. $\mathbb{R}^3$}
If $B$ is a solid in $\mathbb{R}^3$, $vol(T_A(B)) = |\det(A)|vol(B)$
This means that the volume of the parallelepiped spanned by $\vec{u}, \vec{v}, \vec{w}$ is given by $V = |\det([\vec{u}, \vec{v}, \vec{w}])|$

\pagraph(Cross Product)
The cross product $\vec{u} \times \vec{v}$ is the vector orthogonal to both $\vec{u}$ and $\vec{v}$ with magnitude $|\vec{u}||\vec{v}|\sin(\theta)$. 
It is only defined on $\mathbb{R}^3$ and is given by:
\begin{align*}
\vec{u} \times \vec{v} = \det \begin{vmatrix}
    \hat{i} & \hat(j) & \hat{k} \\
    u_1 & u_2 & u_3 \\
    v_1 & v_2 & v_3
\end{vmatrix} = \hat{i} \begin{vmatrix}
    a_2 & a_3 \\
    b_2 & b_3
\end{vmatrix} - \hat{j} \begin{vmatrix}
    a_1 & a_3 \\
    b_1 & b_3
\end{vmatrix} + \hat{k} \begin{vmatrix}
    a_1 & a_2 \\
    b_1 & b_2
\end{vmatrix}
\end{align*}

